\lecture{7.7}{Targeted Minimum Loss-Based Estimation}
\label{lec:7-7}

\subsection{Compatible Plug-In Estimators}

A compatible plug-in estimator is an estimator that can be written as a parameter
mapping applied to an estimated distribution. For example, the g-computation
estimator
\[
\psi_{n,G}(1)=\frac{1}{n}\sum_{i=1}^n\Qbar_n(1,W_i)
\]
is a plug-in estimator: it applies the counterfactual mean mapping to an estimated
outcome regression and empirical covariate distribution.

IPTW and AIPTW are not generally plug-in estimators in this sense. They can produce
values outside natural parameter bounds. For example, for a binary outcome, a
counterfactual mean should lie in $[0,1]$, but a weighted residual estimator can in
finite samples fall outside that interval.

\subsection{Motivation for TMLE}

Targeted minimum loss-based estimation constructs a plug-in estimator with
influence-function-based bias correction. The idea is to start with an initial outcome
regression $\Qbar_n$, then update it in a targeted way so that the empirical mean of
the efficient influence function is approximately zero.

For $\psi(1)$, the relevant residual equation is
\[
0=
\frac{1}{n}\sum_{i=1}^n
\frac{\ind(A_i=1)}{g_{n,1}(W_i)}
\{Y_i-\Qbar_n^*(1,W_i)\}.
\]
If the updated regression $\Qbar_n^*$ satisfies this equation and remains a valid
conditional mean, then the plug-in estimator
\[
\psi_{n,\TMLE}(1)=\frac{1}{n}\sum_{i=1}^n\Qbar_n^*(1,W_i)
\]
inherits the same first-order bias correction as AIPTW while preserving parameter
constraints.

\subsection{Initial Outcome Regression}

TMLE begins by estimating $\Qbar(A,W)$ and $g_1(W)$. These can be estimated using
parametric models or flexible learners such as Super Learner. For a bounded or binary
outcome, the initial outcome regression should produce predictions inside $(0,1)$.

\subsection{Targeting Step}

For simplicity, consider a binary outcome and estimation of $\psi(1)$. Let
\[
K_i=\logit\{\Qbar_n(1,W_i)\}
\]
be the offset and define the clever covariate
\[
H_1(A_i,W_i)=\frac{\ind(A_i=1)}{g_{n,1}(W_i)}.
\]
Fit a logistic regression with outcome $Y$, offset $K_i$, and single covariate
$H_1(A_i,W_i)$, often using only observations with $A_i=1$ for this one-treatment
update. The fitted coefficient is denoted $\hat\epsilon$. Logistic fluctuation is
directly natural for binary outcomes or outcomes scaled to $[0,1]$; for
continuous unbounded outcomes, TMLE should use an appropriate loss and submodel
or transform the outcome before using a bounded-loss formulation.

\subsection{Clever Covariate}

The clever covariate is called ``clever'' because it is chosen from the efficient
influence function. For treatment level $1$,
\[
H_1(A,W)=\frac{\ind(A=1)}{g_1(W)}.
\]
It weights the outcome-regression fluctuation so that the updated regression solves
the empirical residual equation associated with $\psi(1)$.

\subsection{Updated Regression}

The updated regression for treatment level $1$ is
\[
\Qbar_n^*(1,w)
=
\expit\left[
\logit\{\Qbar_n(1,w)\}
+
\hat\epsilon\frac{1}{g_{n,1}(w)}
\right].
\]
More generally, to target both $\psi(1)$ and $\psi(0)$ simultaneously, define
\[
H_1(A,W)=\frac{A}{g_{n,1}(W)},\qquad
H_0(A,W)=\frac{1-A}{g_{n,0}(W)}.
\]
Fit a logistic fluctuation model with offset $\logit\{\Qbar_n(A,W)\}$ and covariates
$H_1$ and $H_0$. If the fitted coefficients are $\hat\epsilon_1$ and
$\hat\epsilon_0$, then
\[
\Qbar_n^*(a,w)
=
\expit\left[
\logit\{\Qbar_n(a,w)\}
+
\hat\epsilon_1\frac{a}{g_{n,1}(w)}
+
\hat\epsilon_0\frac{1-a}{g_{n,0}(w)}
\right].
\]

\subsection{TMLE Estimate}

The TMLE counterfactual means are plug-in estimates:
\[
\psi_{n,\TMLE}(1)
=
\frac{1}{n}\sum_{i=1}^n\Qbar_n^*(1,W_i),
\]
\[
\psi_{n,\TMLE}(0)
=
\frac{1}{n}\sum_{i=1}^n\Qbar_n^*(0,W_i),
\]
and
\[
\widehat{\ATE}_{\TMLE}
=
\psi_{n,\TMLE}(1)-\psi_{n,\TMLE}(0).
\]

\begin{implementationbox}
Conceptual TMLE code for a binary outcome:
\begin{lstlisting}[language=R]
H1 <- A / g_1W
H0 <- (1 - A) / g_0W
K  <- qlogis(Qbar_AW)

target_fit <- glm(Y ~ -1 + offset(K) + H1 + H0,
                  family = binomial(),
                  start = c(0, 0))

eps1 <- coef(target_fit)["H1"]
eps0 <- coef(target_fit)["H0"]

Qstar_1W <- plogis(qlogis(Qbar_1W) + eps1 / g_1W)
Qstar_0W <- plogis(qlogis(Qbar_0W) + eps0 / g_0W)

tmle_ate <- mean(Qstar_1W - Qstar_0W)
\end{lstlisting}
\end{implementationbox}

\subsection{TMLE vs AIPTW}

AIPTW and TMLE have the same efficient influence function under standard conditions.
Both are doubly robust for consistency. The key difference is that TMLE is a targeted
plug-in estimator. For bounded outcomes, it respects the natural bounds of the
parameter. It can also have better finite-sample behavior in settings with practical
positivity concerns.

\begin{keyidea}
TMLE updates the initial outcome regression just enough to solve the efficient
influence-function estimating equation, then estimates the causal parameter by
plugging in the updated regression.
\end{keyidea}
